
\documentclass[letter, 11pt]{article}

    %%%%%%%%%%%%%%%%%%%%%%%%%%%%%%
    \oddsidemargin  0.0in
    \evensidemargin 0.0in
    \textwidth      6.8in
    \headheight     0.0in
    \topmargin      0.0in
    \textheight=9.0in
    \parindent=0in
    %%%%%%%%%%%%%%%%%%%%%%%%%%%%%%


\renewcommand{\familydefault}{\sfdefault} % Police 

\usepackage{amsthm,amsfonts,amssymb,amsmath, xcolor, enumitem}
\usepackage{graphicx}   %To include graphics
\usepackage[small,compact]{titlesec}

\usepackage[utf8]{inputenc}

\usepackage[french]{babel}

\usepackage{fancyhdr}
\pagestyle{fancy}

\newcommand{\var}{\mathbb{V}\text{ar}} %Symboles programmés pour simplifier le texte



\usepackage{version} % Pour faire un texte troué pour les étudiants
%\includeversion{etudiant}\excludeversion{professeur} %version etudiant
\includeversion{professeur}\excludeversion{etudiants} %version prof
\newcommand{\etudiant}[1]{\begin{etudiants}#1 \end{etudiants}}
\newcommand{\prof}[1]{\begin{professeur}\color{blue} {\bf Solution :} #1 \end{professeur}}


\fancyhead[L]{STT-2100 \\ Régression}
\fancyhead[R]{A.-S. Charest, E. Reny-Nolin, A.-A. Vallée \\ Université Laval}

\newcommand{\N}{\mathcal{N}}
\newcommand{\bo}[1]{\boldsymbol#1}
\begin{document}

\vspace*{0.01 cm}
\begin{center}
{\Large Exercices - Chapitre 4}
\end{center}




\begin{enumerate}[leftmargin=*, wide, label=\bf Exercice \arabic*:]

%%%%%%%%%%%%
% Question 1
\item Soit la matrice $\bold{A}=\left[\begin{array}{cc}3&1\\4&2\end{array}\right]$ et le vecteur $\bold{v}'=[v_1\;v_2]$.

\bigskip
    \begin{itemize}
    \item[(a)] Déterminez $\dfrac{\partial \bold{v'A}}{\partial \bold{v}}$.
    \item[(b)] Déterminez $\dfrac{\partial \bold{v'Av}}{\partial \bold{v}}$.
    \end{itemize}

\prof{ 
    \begin{itemize}
    \item[(a)]  
    \begin{align*}
        \bf{v'A}
        &= \begin{pmatrix}
            v_1 & v_2
        \end{pmatrix}
        \begin{pmatrix}
            3 & 1 \\[4mm]
            4 & 2
        \end{pmatrix}
        = \begin{pmatrix}
            3v_1 + 4v_2 & v_1 + 2v_2
        \end{pmatrix}\\
        %
        \dfrac{\partial \bf{v'A}}{\partial \bf{v}}
        &= \begin{pmatrix}
            \dfrac{\partial \bf{v'A}}{\partial v_1} \\[4mm]
            \dfrac{\partial \bf{v'A}}{\partial v_2}
        \end{pmatrix}
        = \begin{pmatrix}
            \dfrac{ 3v_1 + 4v_2}{\partial v_1} & \dfrac{ v_1 + 2v_2}{\partial v_1} \\[4mm]
            \dfrac{ 3v_1 + 4v_2}{\partial v_2} & \dfrac{ v_1 + 2v_2}{\partial v_2} 
        \end{pmatrix}
        = \begin{pmatrix}
            3 & 1 \\[4mm]
            4 & 2
        \end{pmatrix}
        = \bold A
    \end{align*}
    \item[(b)] \begin{align*}
        \bf{v'Av}
        &= \begin{pmatrix}
            v_1 & v_2
        \end{pmatrix}
        \begin{pmatrix}
            3 & 1 \\[4mm]
            4 & 2
        \end{pmatrix}
        \begin{pmatrix}
            v_1 \\[4mm] v_2
        \end{pmatrix}
        = (3v_1 + 4v_2)v_1 + (v_1 + 2v_2)v_2\\
        %
        \dfrac{\partial \bf{v'Av}}{\partial \bf{v}}
        &= \begin{pmatrix}
            \dfrac{\partial \bf{v'Av}}{\partial v_1} \\[4mm]
            \dfrac{\partial \bf{v'Av}}{\partial v_2}
        \end{pmatrix}
        = \begin{pmatrix}
            6v_1+5v_2 \\[4mm]
            5v_1 + 4v_2
        \end{pmatrix}
        = \begin{pmatrix}
            6 & 5 \\[4mm]
            5 & 4
        \end{pmatrix}
        {\bold v}
        = (\bf{A}+\bf{A}')\bf{v}
    \end{align*}
    \end{itemize}}

%%%%%%%%%%%%%%%%%%%%%%%%%%%%%%%%%%%%%%%%%%%%%%%%%%%%%%%%%%%%%%%%%%%%%%%%%%%%%%%%%%%%%%%%%%
\item Montrez que la matrice de projetion $\bold P_{\bold X} = \bold X \left( \bold X^\top \bold X\right)^{-1} \bold X^\top$
    \begin{itemize}
    \item[(a)] est idempotente
    \item[(b)] est symétrique 
    \item[(c)] a une trace égale à $p' = p+1$. (Utilisez le fait que  tr$(AB)$ = tr$(BA)$ si les dimensions le permettent.)
    \end{itemize}

\prof{
 \begin{itemize}
    \item[(a)]$\bold H\bold H = \bold X(\bold X^\top\bold X)^{-1}\bold X^\top\bold X(\bold X^\top\bold X)^{-1}\bold X^\top=\bold X(\bold X^\top\bold X)^{-1}\bold X^\top=\bold H$

    \item[(b)]$\bold H$ est une matrice symétrique si et seulement si $\bold H^\top=\bold H$. On voit  que c'est le cas:
    $$\bold H^\top = \left(\bold X(\bold X^\top\bold X)^{-1}\bold X^\top\right)^\top=(\bold X^\top)^\top\left((\bold X^\top\bold X)^{-1}\right)^\top\bold X^\top
    =\bold X(\bold X^\top\bold X)^{-1}\bold X^\top = \bold H$$
    \item[(c)]On utilise le fait que  tr$(AB)$ = tr$(BA)$ avec $A=\bold X$ et $B=(\bold X'\bold X)^{-1}\bold X'$. Alors,
    $$\mbox{tr}(\bold H) = \mbox{tr}(\bold X(\bold X^\top\bold X)^{-1}\bold X^\top)=\mbox{tr}((\bold X^\top\bold X)^{-1}\bold X^\top\bold X)=\mbox{tr}({\bold{I}_{p'}})=p'.$$
    \end{itemize}
}

%%%%%%%%%%%%%%%%%%%%%%%%%%%%%%%%%%%%%%%%%%%%%%%%%%%%%%%%%%%%%%%%%%%%%%%%%%%%%%%%%%%%%%%%%%
\item On a \begin{eqnarray*}
\bo{X} = \begin{pmatrix}
X_1 \\[4mm]
X_2 \\[4mm]
\vdots \\[4mm]
X_n \\[4mm]    
\end{pmatrix} \sim \N_n \left(\bo\mu, \bo\Sigma\right).
\end{eqnarray*}
\begin{itemize}
    \item[(a)] Supposons que $X_i \sim \N \left(0, 1\right)$ pour $i = 1, \dots, n$ et que les $n$ variables aléatoires sont indépendantes. Écrivez $\bo\mu$ et $\bo\Sigma$ dans ce cas.
    \item[(b)] Supposons que tous les $X_i$ sont des variables aléatoires normales homoscédastiques, de même moyenne et indépendantes. Écrivez $\bo\mu$ et $\bo\Sigma$ dans ce cas.
    \item[(c)] Supposons que tous les $X_i$ sont des variables aléatoires normales hétéroscédastiques, de même moyenne et indépendantes. Écrivez $\bo\mu$ et $\bo\Sigma$ dans ce cas.
    \item[(d)] Supposons que $n=4$, que tous les $X_i$ sont des variables aléatoires normales, que la moyenne de $X_1,\, X_2$ et $X_3$ est la même, que la variance de $X_1,\, X_2$ et $X_3$ est la même, que la covariance entre $X_1$ et $X_2$ vaut 4, que la covariance entre $X_3$ et $X_4$ vaut 5 et que les autres variables sont indépendantes. Écrivez $\bo\mu$ et $\bo\Sigma$ dans ce cas.\textbf{}
\end{itemize}

\prof{
\begin{itemize}
    \item[(a)] \begin{align*}
        \bo\mu
        &= \begin{pmatrix}
            0 \\[4mm]
            \vdots \\[4mm]
            0
        \end{pmatrix}\\
        \bo\Sigma
        &= \begin{pmatrix}
            1 & 0 & \cdots & 0 & 0 \\[4mm]
            0 & 1 & \cdots & 0 & 0 \\[4mm]
            \vdots & & \ddots & & \vdots \\[4mm]
            0 & 0 & \cdots & 0 & 1
        \end{pmatrix} = \mathcal{I}_n
    \end{align*}
    \item[(b)] \begin{align*}
        \bo\mu
        &= \begin{pmatrix}
            \mu \\[4mm]
            \vdots \\[4mm]
            \mu
        \end{pmatrix}\\
        \bo\Sigma
        &=  \begin{pmatrix}
            \sigma^2 & 0 & \cdots & 0 & 0 \\[4mm]
            0 & \sigma^2 & \cdots & 0 & 0 \\[4mm]
            \vdots & & \ddots & & \vdots \\[4mm]
            0 & 0 & \cdots & 0 & \sigma^2
        \end{pmatrix} = \sigma^2\mathcal{I}_n
    \end{align*}
     \item[(c)] \begin{align*}
        \bo\mu
        &= \begin{pmatrix}
            \mu \\[4mm]
            \vdots \\[4mm]
            \mu
        \end{pmatrix}\\
        \bo\Sigma
        &=  \begin{pmatrix}
            \sigma_1^2 & 0 & \cdots & 0 & 0 \\[4mm]
            0 & \sigma_2^2 & \cdots & 0 & 0 \\[4mm]
            \vdots & & \ddots & & \vdots \\[4mm]
            0 & 0 & \cdots & 0 & \sigma_n^2
        \end{pmatrix} 
    \end{align*}
    \item[(d)] \begin{align*}
        \bo\mu
        &= \begin{pmatrix}
            \mu_1 \\[4mm]
            \mu_1 \\[4mm]
            \mu_1 \\[4mm]
            \mu_4
        \end{pmatrix}\\
        \bo\Sigma
        &=  \begin{pmatrix}
            \sigma_1^2 & 4 & 0 & 0 \\[4mm]
            4 & \sigma_1^2 & 0 & 0 \\[4mm]
            0 & 0 & \sigma_1^2 & 5 \\[4mm]
            0 & 0 & 5 & \sigma_4^2
        \end{pmatrix} 
    \end{align*}
\end{itemize}
}

\end{enumerate}

\end{document}

